%%%%%%%%%%%%%%%%%%%%%%%%%%%%%%%%%%%%%%%%%%%%%%%%%%%%%%%%%%%%%%%%%%%%%%%%%%%%%%%%%%%%%%%%%%%%%%%%%%%%%%
%
%   Filename    : chapter_3.tex 
%
%   Description : This file will contain your Theoretical Framework.
%                 
%%%%%%%%%%%%%%%%%%%%%%%%%%%%%%%%%%%%%%%%%%%%%%%%%%%%%%%%%%%%%%%%%%%%%%%%%%%%%%%%%%%%%%%%%%%%%%%%%%%%%%

\chapter{Computer Graphics and Rendering Concepts}
\label{sec:theoframework}
%Please refer to the following resources regarding Theoretical Framework: 

The modernization of AnitoPlume involves migrating the rendering pipeline from a traditional to a modern graphics API architecture. 
This migration introduces a fundamentally different graphics architecture characterized by explicit resource management, reduced driver overhead, and improved utilization of modern multicore processors.
These architectural changes are beneficial for simulators such as AnitoPlume, which require continuous updates of thousands of dynamic particle billboards.

\section{Graphics Rendering Architectures}
Traditional graphics API architectures are designed to abstract low-level graphics operations from the programmer by implicitly managing graphics state, resource synchronization, and command validation within the graphics driver.
Throughout this study, these are referred to as implicit graphics APIs, whereas architectures that expose these responsibilities directly to the application are referred to as explicit graphics APIs.

\subsection{Overview of Graphics API Architectures}
Implicit graphics APIs such as OpenGL are based on a state-based programming model where the application modifies global rendering states that are maintained by the graphics implementation \cite{opengl46spec2022}.
Implicit graphics API architecture relies on a thick driver layer that abstracts graphics pipeline management from the application (Fig.~\ref{fig:api_architecture_comparison}).
Because the driver is responsible for maintaining graphics state, performing runtime validation, tracking resource usage, and managing command execution, significant CPU overhead is introduced during rendering (Ping et al., 2026).

Explicit graphics architectures expose synchronization, resource management, and command submission responsibilities directly to applications \cite{vulkan14spec2024}.
Rather than relying on implicit driver-managed state changes, explicit graphics architectures require the application to explicitly define and manage the configuration of the graphics pipeline.
Applications explicitly define pipeline states, resource allocation, synchronization, and memory management as seen on figure \ref{fig:api_architecture_comparison}.

\begin{table}[htbp]
\centering
\small % Reduces font size to save vertical space
\renewcommand{\arraystretch}{1.2} % Modest padding instead of 1.5
\begin{tabular}{|p{4.2cm}|p{4.2cm}|p{4.2cm}|}
\hline
\textbf{OpenGL} & \textbf{Vulkan} & \textbf{DirectX 12} \\ 
\hline
API originally created for graphics workstations; uses split memory. & Match for modern platforms (mobile/unified memory/tiled rendering). & Modern API for Windows/Xbox; scales across unified platforms. \\ 
\hline
Driver handles state validation/error tracking, limiting performance. & Explicit API gives application direct, predictable GPU control. & Shifting validation to driver, giving developers explicit resource control. \\ 
\hline
Obsolete single-threaded model; no parallel command generation. & Built for multi-core/threads; parallel command buffers. & Multi-core optimized; pre-computes command lists on any CPU core. \\ 
\hline
Complex choices; syntax evolved over 20+ years. & No legacy requirements; simpler API design and smaller specification. & Replaces legacy state objects with unified Pipeline State Objects (PSO). \\ 
\hline
In-driver GLSL compiler; requires shipping raw shader source. & SPIR-V target enables front-end language flexibility and reliability. & Compiles HLSL to intermediate DXIL/DXBC to reduce driver workload. \\ 
\hline
High vendor-specific implementation variability. & Simpler API and testing improve cross-vendor compatibility. & Strict Feature Levels (e.g., DX12 Ultimate) ensure vendor consistency. \\ 
\hline
\end{tabular}
\caption[Graphics API comparison: OpenGL vs. Vulkan vs. DirectX 12]{Graphics API comparison: OpenGL vs. Vulkan vs. DirectX 12 from \citeauthor{shiraef2016exploratory} and \citeauthor{zoraja2023vanityx}}\label{tab:graphics_api_comparison}
\label{tab:dif_API}
\end{table}

Table \ref{tab:dif_API} summarizes the key architectural and performance differences between the OpenGL, Vulkan, and DirectX 12 graphics APIs. 
It contrasts the single-threaded design of OpenGL with the modern, multi-threaded, and explicitly controlled paradigms introduced by Vulkan and DirectX 12.

The limitations of the traditional graphics API architecture become particularly significant in computationally intensive scientific visualization systems such as AnitoPlume. 
In addition to numerically integrating volcanic plume dynamics including gravity, buoyancy, drag, and wind forces for thousands of particles each simulation step, the CPU must also handle graphics driver validation and state management \cite{del2026interactive}. 
This added overhead limits rendering performance, motivating the use of lower-overhead graphics APIs to reduce CPU-side workload for applications that require high-throughput rendering. 
As rendering workloads become more complex, these driver-managed operations can limit the application's ability to optimize resource usage, creating CPU-side bottlenecks that reduce the rate at which work can be submitted to the GPU \cite{ping2026modern}.

\begin{figure}[H]
    \centering
    \makebox[\textwidth][c]{
        \includegraphics[width=1.2\textwidth]{Implicit_Vs_Explicit.png}
    }
    % \protect is mandatory when using apacite inside captions
    \caption{Comparison of implicit graphics API architecture (left) with explicit graphics API architecture (right).}
    \label{fig:api_architecture_comparison}
\end{figure}


\subsection{Performance Optimization opportunities}
Explicit graphics architectures reduce driver-managed serialization by shifting the responsibility to the application.
Tasks such as pipeline configuration, resource allocation, and command preparation can be performed explicitly by the application before or outside the main rendering loop. During runtime, the system only executes predefined instructions without modifying or accessing the underlying data.
This shift allows command generation to be safely distributed across multiple CPU threads \cite{ping2026modern}.
The lightweight driver architecture introduced by explicit graphics APIs minimizes runtime validation and hidden state management, reducing CPU overhead while providing developers with greater control over rendering operations. By enabling application-driven optimization and more efficient utilization of hardware resources, explicit graphics API architectures are better suited for performance-critical rendering workloads \cite{IoannidisC.2020MRFC}.

\section{GPU Parallel Computing}
\subsection{General Purpose GPU Computing}
The increasing computational demands of scientific visualization and physically based simulation have driven the adoption of Graphics Processing Units (GPUs) as general-purpose parallel computing devices.
Unlike Central Processing Units (CPUs), which consist of a relatively small number of complex processing cores optimized for sequential execution, GPUs consist of many smaller cores organized into parallel execution units designed to maximize throughput for highly parallel workloads.
As illustrated in Figure \ref{fig:cpuVgpu}, this architecture enables GPUs to efficiently process workloads exhibiting high degrees of data parallelism, where identical computations are repeatedly applied across large collections of independent data elements \cite{OwensJohnD.2008GC}.

\begin{figure}[H]
    \centering
    \makebox[\textwidth][c]{
        \includegraphics[width=0.9\textwidth]{cpuVgpu.png}
    }
    % \protect is mandatory when using apacite inside captions
    \caption{CPU vs. GPU key differences. From \protect\cite{TecAdmin2025}.}
    \label{fig:cpuVgpu}
\end{figure}

The emergence of general purpose computing on GPUs transformed graphics hardware from fixed-function rendering accelerators into programmable parallel processors capable of executing scientific and engineering applications.

\subsection{GPU Programming}
Early GPU computing approaches utilized graphics APIs to express numerical algorithms as rendering operations.
This paradigm evolved with the introduction of programming models such as NVIDIA's CUDA, which exposed the GPU as a general parallel computing platform through familiar programming abstractions while providing direct access to the device's computational resources \cite{LuebkeD.2008CSpp}.

\citeauthor{HarrisMark2004Mcct} describes GPU programming as a process of mapping computational problems onto massively parallel hardware by identifying operations that can be executed independently across many processing elements.
Rather than accelerating a single complex computation, GPU applications achieve performance by decomposing a problem into numerous independent tasks that execute simultaneously. Consequently, algorithms exhibiting minimal inter-element dependencies and regular memory access patterns are generally well suited for GPU execution \cite{HarrisMark2004Mcct}.
Compute shaders organize execution into groups of parallel invocations that simultaneously execute the same program on different portions of the input data.
Unlike vertex or fragment shaders, compute shaders are not directly associated with rendering primitives, allowing applications to perform arbitrary numerical computations while explicitly managing memory resources and synchronization between execution stages \cite{vulkan14spec2024}.

Because many scientific simulations consist of repeatedly updating large collections of independent entities, GPU computing has become widely adopted across computational physics, fluid dynamics, molecular simulation, image processing, and real-time visualization. \citeauthor{OwensJohnD.2008GC} identify data-parallel simulation as one of the primary application domains benefiting from GPU architectures due to their ability to execute thousands of lightweight computational threads concurrently. As GPU hardware and programming models continue to evolve, massively parallel computing has become a fundamental approach for accelerating computationally intensive scientific applications while maintaining interactive performance \cite{OwensJohnD.2008GC}.

\subsection{CUDA Programming}
The Compute Unified Device Architecture (CUDA), introduced by NVIDIA in 2006, enables any computational workload to use the capability of GPUs independent of graphics APIs. CUDA is implemented in systems that include both CPUs and GPUs, referred to as the host and device respectively. The host and device have their own respective memory, and can execute code independently. The increased performance of CUDA applications happens by maximizing the utilization of both CPU and GPU resources.

\subsubsection{SIMT and Tile Model}
A kernel is a function executed by multiple GPU threads simultaneously. Threads are organized into thread blocks, and thread blocks are further organized into a grid. Threads within a grid know its location within its block and the location of the block within its grid. Thus each thread has a unique identity among all other threads, allowing programs to determine what data or operations a thread is responsible for. More importantly, threads within a thread block all have access to the on-chip shared memory, which can be used for exchanging information between each other efficiently. 

CUDA employs the Single Instruction, Multiple Threads (SIMT) execution model, in which multiple GPU threads execute the same instruction sequence concurrently while operating on different data elements. Threads are organized into groups of 32 called warps, allowing the GPU to efficiently process large numbers of similar computations in parallel (Figure~\ref{fig:SIMT}).

\begin{figure}[H]
    \centering
    \makebox[\textwidth][c]{
        \includegraphics[width=0.9\textwidth]{SIMT.png}
    }
    % \protect is mandatory when using apacite inside captions
    \caption{Programmer’s view in the SIMT and tile programming models. In SIMT, the programmer writes per-thread code and controls how each thread accesses data. In tile programming, the programmer writes per-block code that operates on tiles; the compiler maps operations to the threads of the block. From \protect\cite{nvidia_tile_simt_figure}.}
    \label{fig:SIMT}
\end{figure}

Tile kernels, which are functions operating on tiles, are launched on a grid of blocks. Tile programming is one way of utilizing CUDA in an application, allowing programs to describe operations on multidimensional collections of data called tiles. Tile kernels work on arrays and tiles, and may create and operate on many tiles of different shapes and data types (Figure~\ref{fig:TileSpace}). 

\begin{figure}[H]
    \centering
    \makebox[\textwidth][c]{
        \includegraphics[width=0.9\textwidth]{Tile.png}
    }
    % \protect is mandatory when using apacite inside captions
    \caption{Tile space and data movement. A two-dimensional array of shape (M, N) is conceptually partitioned into a grid of tiles with shape (tm, tn). A load at tile-space index (i, j) returns the corresponding tile. At the array boundary, elements that fall outside the array can be filled with zeros. A store writes a tile back to the array at a given tile-space index. From \protect\cite{nvidia_tile_space_data_movement_figure}.}
    \label{fig:TileSpace}
\end{figure}

Figure \ref{fig:Code} demonstrates how a tile kernel can be used to perform parallel operations on an array. The kernel calculates the sum of multiple tiles stored in an input array and produces a single output tile. The input array is partitioned into tiles containing eight floating-point values. The kernel creates a view of the input array where each tile can be accessed using a tile index. An accumulator tile is initialized with zero values, and the kernel iterates through all input tiles, loading each tile and performing element-wise addition with the accumulator. After all tiles have been processed, the accumulated result is stored as the output tile. Although the example performs a reduction operation, the same tile-based approach can be applied to simulation workloads where large arrays must be updated repeatedly. Instead of summing tiles, each tile can represent a section of a simulation domain, such as a group of cells in a two-dimensional grid. A tile kernel can load a section of the array, perform computations on the values within that region, and store the updated tile back into memory.

\begin{figure}[H]
    \centering
    \makebox[\textwidth][c]{
        \includegraphics[width=0.9\textwidth]{Code.png}
    }
    % \protect is mandatory when using apacite inside captions
    \caption{CUDA tile kernel execution for summing multiple input tiles into an output tile. From \protect\cite{nvidia_cuda_programming_guide_13_3_tile_sum}.}
    \label{fig:Code}
\end{figure}

\subsubsection{CUDA Memory Types}

CUDA provides a hierarchical memory system designed to optimize data access during parallel execution. For tile-based computation, global memory and shared memory are particularly important. Global memory stores large datasets accessible across the GPU, while shared memory provides fast temporary storage for threads within a block. By loading data into shared memory and processing it locally, tile-based approaches can reduce memory access overhead and improve computational efficiency. 

\begin{table}[H]
\centering
\caption{Global vs Shared Memory in CUDA}
\label{tab:cuda_memory_types}
\begin{tabular}{p{2.2cm} p{2cm} p{2cm} p{2cm} p{4cm}}
\toprule
\textbf{Memory Type} & \textbf{Scope} & \textbf{Lifetime} & \textbf{Location} & \textbf{Purpose} \\
\midrule
Global Memory & Grid & Application & Device & Primary memory space for storing data used by threads \\
\midrule
Shared Memory & Block & Kernel & SM & Memory space accessible by all threads within a thread block \\
\bottomrule
\end{tabular}
\end{table}

There is no guarantee of scheduling between thread blocks, so a thread block cannot rely on results from other thread blocks, as they may not be able to be scheduled until that thread block has completed. Synchronization ensures that concurrent threads execute correctly when accessing shared resources. In CUDA, threads within the same thread block can synchronize their execution and coordinate access to shared memory, while thread blocks execute independently without guaranteed scheduling order. CUDA provides mechanisms such as barrier synchronization and atomic operations to prevent race conditions during parallel execution.

However, CUDA programs must also guarantee synchronization within the host code. Mutexes, locks, and semaphores are commonly used in CPU-based multithreaded systems to control access to shared resources and coordinate concurrent tasks. These synchronization concepts are used in GPU-accelerated rendering systems where CPU-side processes and GPU computations must coordinate efficiently. 

\subsection{Application of Compute Shaders in Parallel Computing}
Data-parallel simulation is one of the primary application domains benefiting from GPU architectures due to their ability to execute thousands of lightweight computational threads concurrently \cite{OwensJohnD.2008GC}. Recent studies demonstrate the use of compute shaders across different simulation and visualization domains, including particle-based simulation, physical deformation models, and large-scale rendering. 

\citeauthor{Va2021RealTimeCloth} further demonstrated the use of compute shaders for iterative physics simulation, where simulation stability and performance depend on repeated updates across discrete time steps. Similarly, \citeauthor{madera2022cpu} applied compute shaders to particle-based sphere simulation by distributing particle updates and collision operations across GPU threads. Their work addressed the increasing computational and memory requirements caused by large particle counts by improving CPU–GPU buffer communication and maintaining a reusable computation buffer. \citeauthor{schutz2021rendering} demonstrated the use of compute shaders for large-scale point cloud rendering. Their approach replaced portions of the traditional graphics pipeline with compute-based processing, allowing independent point data to be processed through GPU parallel execution. Although these works address different applications, they illustrate how compute shaders exploit data parallelism to improve throughput for workloads consisting of many independent elements.

\begin{table}[H]
\centering
\caption{Application of compute shaders for optimization.}
\vspace{0.2cm}
\renewcommand{\arraystretch}{1.5}
% Note: Requires \usepackage{tabularx} in your preamble
\begin{tabularx}{\textwidth}{|>{\hsize=0.6\hsize}X|>{\hsize=0.8\hsize}X|>{\hsize=1.3\hsize}X|>{\hsize=1.3\hsize}X|}
\hline
\textbf{Study} & \textbf{Domain} & \textbf{Compute Shader Application} & \textbf{Results} \\ \hline
Va et al. (2021) & Real-time cloth simulation & Implemented a GPU-based parallel cloth simulation using compute shader kernels for mass--spring calculations and constraint enforcement & GPU parallel implementation achieved faster performance than CPU-based implementation while maintaining stable cloth behavior using constraint-based methods. \\ \hline
Madera-Ramírez et al. (2021) & Particle Simulation & Used compute shaders to parallelize particle updates, collision operations, and GPU buffer communication & Achieved up to $100\times$ speedup over the sequential implementation and enabled processing of larger particle counts. \\ \hline
Schütz et al. (2021) & Point cloud rendering & Used compute shaders to perform point processing and rendering operations through GPU-based parallel execution & Demonstrated high-throughput rendering of hundreds of millions of points with improved performance over previous approaches. \\ \hline
\end{tabularx}
\end{table}

Similar characteristics are present in particle-based plume simulation, where each plume element maintains independent state variables that must be repeatedly updated during each simulation step. Therefore, compute shaders provide a suitable computational model for accelerating both simulation updates and visualization of large numbers of dynamic plume elements. 

\section{Lagrangian and Eulerian Models for Volumetric Plume Dynamics}

\subsection{Lagrangian Formulation}
AnitoPlume utilizes a Lagrangian formulation for representing volcanic plume evolution, where the plume is modeled as a collection of discrete volumetric elements whose physical states are updated over time.
Unlike Eulerian approaches that evaluate flow properties within fixed spatial domains, Lagrangian approaches track the movement and evolution of individual elements within the flow field \cite{EspostiOngaro2020Validation}.

The plume formulation presented by \citeA{lastic2022interactive} and implemented in AnitoPlume \cite{del2026interactive} represents plume dynamics through discrete elements with associated physical properties.
Each plume element maintains state variables such as position, velocity, mass, density, and volume. During each simulation step, these variables are updated based on governing physical equations that consider forces including gravity, buoyancy, atmospheric drag, and wind effects \cite{lastic2022interactive,del2026interactive,gmd-9-431-2016}.

The structure of the Lagrangian formulation exhibits data-parallel characteristics because individual plume elements can be updated primarily using their local state information and shared simulation parameters.
Consequently, operations such as force evaluation and numerical integration can be distributed across multiple processing units, provided that synchronization requirements are properly managed \cite{lastic2022interactive}.

However, Lagrangian particle-based methods also present limitations when applied to large-scale geophysical phenomena. 
The computational cost associated with representing the large number of particles required for realistic volcanic processes has limited the widespread use of fully Lagrangian approaches in volcanological simulations \cite{EspostiOngaro2020Validation}. 
These limitations motivate the consideration of alternative formulations, model refinements, or modifications to the underlying mathematical representation depending on the accuracy and computational requirements of the simulation. 
Consequently, the selection of governing equations and the formulation of the plume model remain important considerations in the continued development of volcanic simulation systems.

Despite these considerations, the data-parallel characteristics of the Lagrangian representation make plume simulation suitable for optimization through modern parallel processing architectures.
Since the computational workload consists of continuously updating large numbers of dynamic elements, improving execution efficiency requires an architecture capable of reducing overhead and efficiently utilizing available hardware resources.
Explicit graphics architectures provide mechanisms that support this objective by enabling greater control over resource management and parallel command execution.


\begin{figure}[H]
    \centering
    \makebox[\textwidth][c]{
        \includegraphics[width=0.85\textwidth]{Langrangian.png}
    }
    % \protect is mandatory when using apacite inside captions
    \caption{Lagrangian approach. From \protect\cite{lastic2022interactive}.}
    \label{fig:lagrangian}
\end{figure}

Figure \ref{fig:lagrangian} illustrates the Lagrangian formulation implemented by \cite{lastic2022interactive}, detailing the physical forces, wind-driven trajectory changes, and volumetric expansion acting on individual discrete plume elements over time.

\subsection{Eulerian Formulation}

In contrast to the Lagrangian approach, the Eulerian formulation describes plume dynamics with reference to fixed points or cells in space, rather than tracking discrete elements as they move through the flow field. 
Under this formulation, the volcanic plume is represented as a continuous field of quantities, such as ash concentration, density, temperature, and velocity, which are computed and updated at fixed grid locations over time as material flows through them \cite{EspostiOngaro2020Validation}. 
Rather than following individual particles or plume elements as they travel through the atmosphere, the Eulerian approach numerically solves governing transport equations directly on a spatial grid that discretizes the simulation domain. 
This formulation underlies many of the numerical geoscientific models used for volcanic hazard forecasting; for instance, models such as FALL3D use Eulerian atmospheric dispersion techniques to compute the transport, dispersion, and eventual deposition of volcanic tephra across a defined spatial domain \cite{folch2009fall3d,del2026interactive}. 
While this grid-based approach is generally well-suited for capturing diffusive processes and producing quantitatively accurate forecasts, it comes at a significant computational cost, particularly as grid resolution increases or the simulation domain grows, since the governing equations must be solved at every grid cell regardless of whether volcanic material is present at that location. 
Consequently, Eulerian volcanic transport models are typically designed for offline scientific analysis and hazard mapping rather than for real-time or interactive visualization \cite{del2026interactive}.

\begin{figure}[H]
    \centering
    \makebox[\textwidth][c]{
        \includegraphics[width=0.6\textwidth]{eulerian.png}
    }
    % \protect is mandatory when using apacite inside captions
    \caption{Eulerian approach. From \protect\cite{folch2020fall3d}.}
    \label{fig:eulerian}
\end{figure}

For the Eulerian Approach, Figure \ref{fig:eulerian} visualizes how it discretizes the simulation space into a 3D staggered mesh of fixed computational cells to evaluate mass transport and flow velocity over time.

\subsection{Comparison of Both Approaches}

\begin{table}[htbp]
\centering
\small
\caption{Summary of key theoretical distinctions between the Lagrangian and Eulerian formulations.}
\label{tab:lagrangian_eulerian_comparison}
\begin{tabular}{l p{5cm} p{5cm}}
\hline
\textbf{Dimension} & \textbf{Lagrangian Formulation} & \textbf{Eulerian Formulation} \\ \hline
\textbf{Frame of reference} & Tracks discrete, moving plume elements. & Evaluates fixed spatial grid points. \\
\textbf{Representation} & Discrete particles with individual states (pos, vel, mass). & Continuous field of quantities (concentration, density). \\
\textbf{Computation} & Scales with particle count; local-state updates. & Scales with grid size; solved at every cell. \\
\textbf{Accuracy} & Approximates diffusive behavior. & High quantitative accuracy for diffusive transport. \\
\textbf{Real-time} & Well-suited; lightweight for moderate counts. & Generally unsuitable; computationally demanding. \\
\textbf{Representative use} & Lastic et al. (2022); AnitoPlume & FALL3D (Folch et al., 2009) \\
\textbf{Limitation} & Cost grows with particle count for visual realism. & Cost grows with grid resolution/domain coverage. \\ \hline
\end{tabular}
\end{table}

AnitoPlume adopts the Lagrangian formulation because it prioritizes real-time interactivity over the grid-based accuracy of Eulerian models. 
This study preserves that formulation as a fixed theoretical foundation, since its scope is limited to the rendering and software architecture layer, not the underlying plume dynamics. 
The Lagrangian model's data-parallel characteristics, individual elements updating largely independently using local state, are what make it a suitable target for the proposed low-level rendering architecture, which is designed to expose fine-grained control over parallel command execution and GPU resource management. 
This theoretical basis justifies why the present study's contributions are directed at the rendering pipeline rather than the plume simulation model itself.
%\begin{itemize}

%\item \url{https://link.springer.com/chapter/10.1007/978-1-4419-1454-5_12}
%\item \url{https://link.springer.com/article/10.1007/s10972-015-9443-2}



%\end{itemize}

\section{Profiling Tools}
\label{sec:profiling_tools}

\subsection{The Role of Profiling in Graphics Benchmarking}
\label{sec:profiling_importance}

Profiling is the process of instrumenting a software system to collect quantitative data about its runtime behavior: execution time, resource utilization, throughput, and synchronization patterns \cite{cheung2024applications}. The purpose is to identify performance bottlenecks and validate optimization efforts.
In graphics development, profiling operates across two domains: the CPU side, where the application records and submits commands, and the GPU side, where those commands are translated into rendering and compute workloads \cite{ping2026modern}.
For scientific visualization applications such as AnitoPlume, where a single frame may involve thousands of particle billboard draws, a terrain pass, and a UI composite, profiling is the only means of determining precisely where frame time is being consumed.

In implicit graphics APIs, profiling is primarily concerned with identifying draw-call overhead generated by heavy driver validation and state tracking \cite{shiraef2016exploratory}.
In explicit graphics APIs, the profiler surfaces a broader set of implementation concerns: the cost of recording command buffers, the impact of pipeline barrier placement on GPU occupancy, the efficiency of descriptor set layouts, whether pipeline state objects are reused effectively, and whether compute-to-render transitions are properly scheduled to avoid stalling the pipeline \cite{ping2026modern,vulkan14spec2024}.
These concerns go beyond the quality of the rendered output, and failure to profile them results in a system that may produce correct visuals but fall short of the hardware's capabilities \cite{IoannidisC.2020MRFC}.

The following subsections survey three profiling tools relevant to explicit API development and benchmarking. These tools differ in scope, ranging from real-time statistical profiling to offline frame debugging to deep hardware performance analysis, and together illustrate the instrumentation dimensions that modern graphics development demands.

\subsection{Tracy Profiler}
\label{sec:tracy}

Tracy is an open-source, cross-platform profiler that works across all GPU hardware vendors and is designed for real-time, frame-based performance analysis \cite{tracy105}.
It requires manual source-code instrumentation: the developer inserts zone-marking macros around blocks of interest, such as the simulation update loop, command buffer recording, and presentation. Tracy records the start time, duration, and thread affinity for each zone across every captured frame, building statistical profiles over time.

On the GPU side, Tracy supports recording timestamp queries directly within command buffers, allowing the profiler to track GPU-side execution activity on the same timeline as CPU zones \cite{tracy105}.
This integration is significant for evaluating parallelization: the developer can observe whether a worker thread finishes updating its assigned smoke layers well before others, indicating load imbalance, or whether GPU execution overlaps with CPU command recording rather than idling while waiting for CPU submission.

Tracy is suited for identifying trends that manifest over multiple frames, such as frame timing spikes, CPU-side bottlenecks in the physics integration pass, and idle gaps caused by thread barriers. It is less suited for per-event GPU-level diagnostics, which require offline inspection of individual frames.

\begin{figure}[H]
    \centering
    \makebox[\textwidth][c]{
        \includegraphics[width=1\textwidth]{Tracy.png}
    }
    % \protect is mandatory when using apacite inside captions
    \caption{Tracy Profiler.}
    \label{fig:tracy}
\end{figure}

\subsection{NVIDIA Nsight Graphics}
\label{sec:nsight}

NVIDIA Nsight Graphics is a GPU profiling and debugging tool that provides access to low-level hardware and performance counters available on NVIDIA GPUs \cite{nsightngrade}.
Unlike generic frame capture tools, Nsight Graphics directly queries hardware performance counters including shader processor occupancy, texture unit utilization, memory bandwidth throughput, L1 and L2 cache hit rates, and warp stalls for each draw or dispatch call within a captured frame.

These metrics provide resolution that a generic frame-capture profiler cannot offer.
For example, identifying that a fragment shaded render pass for volcanic plume billboards is bound by texture memory latency rather than arithmetic throughput suggests specific optimization steps: reducing texture resolution, packing multiple data channels into a single texture, or altering the filtering mode. Each step can be validated against the same hardware counter data.

Because Nsight Graphics is restricted to NVIDIA hardware, it is most appropriate for detailed single-frame profiling on NVIDIA workstations rather than continuous monitoring across diverse hardware. For research groups that deploy NVIDIA GPU workstations, Nsight provides the deepest constituent view of hardware performance available.

\begin{figure}[H]
    \centering
    \makebox[\textwidth][c]{
        \includegraphics[width=1\textwidth]{Nsight.jpg}
    }
    % \protect is mandatory when using apacite inside captions
    \caption{NVIDIA Nsight Graphics.}
    \label{fig:nsight}
\end{figure}

\subsection{RenderDoc}
\label{sec:renderdoc}

RenderDoc is a free, open-source graphics debugging tool that supports Vulkan, DirectX~12, DirectX~11, and OpenGL across GPUs from NVIDIA, AMD, and Intel, making it the most broadly applicable tool surveyed here \cite{renderdoc2023}.
It captures the full sequence of graphics API calls for a single user-specified frame, serializing every draw call, compute dispatch, render pass transition, and memory barrier into an offline trace.

The capture workflow is straightforward: the application is launched through RenderDoc's process injector, a frame is captured via the overlay or hotkey, and the resulting trace is opened in the inspection interface. That interface provides the following:

\begin{itemize}
    \item \textbf{Event Browser.} The time-ordered list of every API call in the frame, allowing the developer to verify that command sequences, barriers, and render pass transitions are executed in the intended order.
    \item \textbf{Pipeline State View.} The complete GPU state, including bound shaders, descriptor sets, vertex and index buffers, rasterizer state, depth-stencil state, and blend state, for any selected event. This enables verification that Pipeline State Objects are correctly constructed and bound.
    \item \textbf{Resource Inspector.} A per-pixel visualization of any intermediate render target, texture, or buffer at any point inside the frame. This is essential for confirming that the terrain pass, plume particle pass, and UI composite pass each produce correct intermediate outputs.
    \item \textbf{GPU Timing.} Per-event GPU-side execution time derived from timestamp queries, identifying which draw calls or compute dispatches are the most expensive within a single frame.
\end{itemize}

\begin{figure}[H]
    \centering
    \makebox[\textwidth][c]{
        \includegraphics[width=1\textwidth]{RenderDoc.jpg}
    }
    % \protect is mandatory when using apacite inside captions
    \caption{RenderDoc.}
    \label{fig:renderdoc}
\end{figure}


\subsection{Summary of Profiling Tools}
\label{sec:profiling_summary}

Each of the profiling tools discussed addresses a distinct layer of graphics performance analysis, and their combined use provides comprehensive instrumentation coverage for projects migrating between graphics APIs. Figures~\ref{fig:tracy}, \ref{fig:nsight}, and \ref{fig:renderdoc} illustrate the user interface of each profiling tool discussed in this section.
Table~\ref{tab:profiling-tools} below contrasts them across five dimensions particularly relevant to this study.

\begin{table}[htbp]
\centering
\small
\renewcommand{\arraystretch}{1.2}
\caption{Comparative summary of the profiling tools discussed.}
\label{tab:profiling-tools}
\begin{tabular}{|p{2.8cm}|p{3.5cm}|p{3.5cm}|p{3.5cm}|}
\hline
\textbf{Dimension} & \textbf{Tracy} & \textbf{NVIDIA Nsight Graphics} & \textbf{RenderDoc} \\
\hline
\textbf{Profiling Mode} & Real-time, continuous across frames & Single-frame, offline & Single-frame, offline \\
\hline
\textbf{Coverage} & CPU zones and GPU timer queries on shared timeline & Per-event hardware counter metrics (occupancy, cache, bandwidth) & Full API call trace, pipeline state, GPU timing \\
\hline
\textbf{GPU Vendor Support} & All (AMD, Intel, NVIDIA) & NVIDIA only & All (AMD, Intel, NVIDIA) \\
\hline
\textbf{API Support} & Any (via vendor-ext timestamp queries) & Vulkan, DirectX~12, DirectX~11, OpenGL & Vulkan, DirectX~12, DirectX~11, OpenGL \\
\hline
\textbf{Primary Use Case} & Identifying CPU-side bottlenecks over multiple frames; validating multithreaded parallelism & Diagnosing hardware-level stalls (cache misses, warp divergence, low occupancy) & Inspecting pipeline state correctness and per-event GPU timing across APIs \\
\hline
\end{tabular}
\end{table}

